%% Supplementary Material S1.
%% Build:  python scripts/make_supplement.py   (writes build/tables, build/fig_*)
%%         cd supplement && pdflatex supplementary_material && pdflatex supplementary_material
%% Anonymous, to match the manuscript under review.
\documentclass[10pt,a4paper]{article}
\usepackage[T1]{fontenc}
\usepackage[utf8]{inputenc}
\usepackage{mathptmx}
\usepackage[a4paper,margin=2.2cm]{geometry}
\usepackage{amsmath,amssymb}
\usepackage{graphicx}
\usepackage{booktabs}
\usepackage{array}
\usepackage[font=small,labelfont=bf,labelsep=period]{caption}
\usepackage{placeins}
\usepackage{xurl}
\usepackage[hidelinks]{hyperref}

\renewcommand{\theequation}{S\arabic{equation}}
\renewcommand{\thesection}{S\arabic{section}}
\renewcommand{\thesubsection}{S\arabic{section}.\arabic{subsection}}
\graphicspath{{build/}}
\setlength{\parskip}{3pt}
\setlength{\parindent}{0pt}

\begin{document}

\begin{center}
{\large\bfseries Supplementary Material S1}\\[6pt]
{\normalsize Two origins of feature redundancy in stroke biosignals: algebraic
construction in transcranial Doppler, shared physiology in quantitative EEG}\\[8pt]
{\small Supplementary methods, Supplementary Tables S1--S11 and Supplementary
Figures S1--S2. The complete numeric results are provided separately as
Supplementary Data S2 (Excel workbook).}
\end{center}

\noindent References to tables, figures and sections without the prefix S refer
to the main text.

\tableofcontents

%=============================================================================
\section{Supplementary methods}
%=============================================================================

The main text describes each analysis and the reasons for it. This section
gives the formulas and settings needed to reproduce them. The analysis code
reproduces every table and figure here from the result files listed in
Supplementary Data S2.

\subsection{Correction for overlapping cross-validation folds}
\label{sec:s-nb}

Every paired comparison uses the $J = 25$ fold-wise differences
$d_1,\dots,d_J$ of the five-times-repeated, five-fold, subject-grouped
evaluation. The folds share training data, so the $d_j$ are not independent,
and the naive variance of their mean, $s^2/J$, is too small. Following Nadeau
and Bengio, the variance of the mean is inflated to
\begin{equation}
\widehat{\operatorname{Var}}(\bar d) \;=\;
\Bigl(\frac{1}{J} + \frac{n_{\mathrm{test}}}{n_{\mathrm{train}}}\Bigr)\, s^2 ,
\qquad s^2 = \frac{1}{J-1}\sum_{j}(d_j-\bar d)^2 .
\label{eq:nb}
\end{equation}
With five-fold partitions each test set holds a fifth of the data, so
$n_{\mathrm{test}}/n_{\mathrm{train}} = 1/4$ in both cohorts (on average 400 of
2000 qEEG trials and 190.2 of 951 TCD insonations per test fold). The same
inflation is applied to the paired significance tests, to the equivalence
tests, and to the power analysis. With $J = 25$ the correction term contributes
$0.25$ against $0.04$ for $1/J$, so the corrected standard error is
$\sqrt{7.25} \approx 2.69$ times the uncorrected one. Because the term
$n_{\mathrm{test}}/n_{\mathrm{train}}$ does not depend on $J$, additional
repeats of the same cross-validation cannot reduce it: only a larger cohort
can.

\subsection{Equivalence tests and their power}

Each comparison of the full and reduced feature spaces is a two one-sided test
(TOST) of $H_0\colon|\mu_d| \ge \Delta$ against $H_1\colon|\mu_d| < \Delta$,
with margin $\Delta = 0.05$ fixed before analysis. With
$\mathrm{SE} = \widehat{\operatorname{Var}}(\bar d)^{1/2}$ from
Eq.~(\ref{eq:nb}),
\begin{equation}
t_{\mathrm{L}} = \frac{\bar d + \Delta}{\mathrm{SE}}, \qquad
t_{\mathrm{U}} = \frac{\bar d - \Delta}{\mathrm{SE}}, \qquad
p_{\mathrm{TOST}} = \max\bigl\{1 - F_{J-1}(t_{\mathrm{L}}),\; F_{J-1}(t_{\mathrm{U}})\bigr\},
\end{equation}
where $F_{\nu}$ is the Student-$t$ distribution function. Equivalently,
equivalence is declared when the $90\%$ interval
$\bar d \pm t_{0.95,\,J-1}\,\mathrm{SE}$ lies inside $(-\Delta, \Delta)$.
Holm adjustment is applied across all 40 tests (five classifiers, two
modalities, four metrics).

Power is evaluated at a true difference of zero, approximating the joint
rejection probability with two noncentral $t$ distributions:
\begin{equation}
1-\beta \;=\; \max\Bigl\{0,\; G_{J-1,\,-\Delta/\mathrm{SE}}\bigl(-t_{0.95,\,J-1}\bigr)
\;-\; G_{J-1,\,+\Delta/\mathrm{SE}}\bigl(t_{0.95,\,J-1}\bigr)\Bigr\},
\end{equation}
where $G_{\nu,\lambda}$ is the noncentral-$t$ distribution function. Power is
zero whenever $\Delta < t_{0.95,\,J-1}\,\mathrm{SE}$: the margin is then
narrower than the half-width the interval would have even for an observed
difference of exactly zero, and no outcome can establish equivalence. The
minimum detectable margin $\delta_{80}$ is the smallest $\Delta$ for which
$1-\beta \ge 0.80$, found by bisection.

\subsection{Threshold sweep}

The reduction chain (correlation screening at $|r| \ge \tau_r$, then
iterative removal of the descriptor with the largest variance inflation factor
while any exceeds $\tau_v$) was recomputed for $\tau_r \in \{0.30, 0.35,
\dots, 0.95\}$ (14 values) and $\tau_v \in \{2.5, 5, 10, 20\}$, on the native
set and the matched core of each modality: $14 \times 4 \times 2 \times 2 =
224$ chains. For each of the 56 threshold pairs and each feature space, the
fraction of candidates retained by qEEG was compared with the fraction retained
by TCD (Supplementary Table~S9; every chain in Supplementary Data S2).

\subsection{Inference on the redundancy gap}

The gap is the difference in mean absolute off-diagonal correlation over the
matched core, TCD minus qEEG. Its interval is a percentile bootstrap with
$2000$ replicates in which \emph{subjects} -- not records -- are resampled with
replacement, independently within each cohort, and every record of a resampled
subject is kept. The supporting permutation check pools the off-diagonal
$|r|$ values of the two matched cores, reassigns them at random to groups of
the observed sizes $2000$ times, and reports the one-sided $p =
(1 + \#\{\text{null} \ge \text{observed}\})/(1 + 2000)$. That null asks
whether a gap this large arises between two arbitrary sets of correlations of
these sizes; it is weaker than the bootstrap and is reported only as support.

\subsection{Redundancy beyond algebraic construction}

Within each native set, descriptors were standardised, and every descriptor
other than the generating set was regressed on that set by ordinary least
squares: the four band powers for qEEG, and peak systolic, end-diastolic and
mean flow velocity for TCD. The residuals' correlation matrix was then
summarised in the same way as the raw one (Supplementary Table~S8).

\subsection{Selection stability}

For a selector applied to $M$ training partitions of a set of $d$
descriptors, let $Z \in \{0,1\}^{M\times d}$ record which descriptors were
retained, $\hat p_f$ the fraction of partitions retaining descriptor $f$, and
$\bar k$ the mean number retained. The Nogueira stability index is
\begin{equation}
\hat\Phi \;=\; 1 - \frac{\frac{1}{d}\sum_{f=1}^{d}\frac{M}{M-1}\,\hat p_f(1-\hat p_f)}
{\frac{\bar k}{d}\Bigl(1-\frac{\bar k}{d}\Bigr)} ,
\end{equation}
which is $1$ for a selector returning the same subset every time and $0$ in
expectation for random subsets of the same size. It was computed over $M = 50$
subject-grouped training partitions. Supervised selectors were run at the
cardinality of the label-agnostic reduced set ($k = 12$ for qEEG, $k = 7$ for
TCD), so that the four selectors are compared at matched subset size.

\subsection{Repeated-measures ANOVA and its decomposition}

Within each modality, fold-wise Macro-F1 was analysed with feature-space
condition (five levels) and classifier (five levels) as within factors and the
25 folds as subjects. Each effect is tested against its interaction with fold;
Greenhouse--Geisser corrected $p$ values and generalised $\eta^2$ are reported
in Table~9. The empirical null permutes the factor labels within each fold
($1000$ permutations). The decomposition restricts the condition factor to the
full and reduced sets, and refits the omnibus test with one supervised selector
removed at a time (Supplementary Table~S11).

\subsection{Supplementary Data S2}

The workbook contains one sheet per result table, including those summarised
here: per-fold metrics for every modality, condition, classifier, repeat and
fold; summaries for all seven metrics; the corrected and uncorrected
equivalence tests and their power; every paired comparison; the chance-level
tests; all 224 threshold-sweep chains; selection frequencies; the stability
indices; the structural and family-block summaries; the partial-correlation
results; the ANOVA and its decomposition; the permutation nulls; the gap
bootstrap; and the channel-pooling, within-subject and calibration results of
the main text. A ``Read me'' sheet lists each sheet with its contents and the
result file it is taken from, and holds formula checks that recompute the
headline equivalence counts from the sheets.

%=============================================================================
\section{Cost of reduction}
%=============================================================================

\FloatBarrier
\input{build/tables/S1}
\input{build/tables/S2}

\begin{figure}[!htbp]
\centering
\renewcommand{\thefigure}{S1}
\includegraphics[width=\linewidth]{fig_S1_equivalence.pdf}
\caption{Equivalence of the full and reduced feature spaces for all 40
comparisons. Thick coloured bars are the $90\%$ intervals with the correction
for overlapping folds (Supplementary Table~S1); thin grey bars above them are
the uncorrected intervals (Supplementary Table~S2). Filled markers are
equivalent after correction and Holm adjustment, open markers are not; the
shaded band is the $\pm0.05$ margin. Positive values are a loss from pruning.
The uncorrected intervals lie inside the margin for most comparisons; once fold
dependence is accounted for, the qEEG intervals remain inside it while the TCD
intervals, centred near zero, become too wide to establish equivalence.}
\label{fig:S1}
\end{figure}

\FloatBarrier
\input{build/tables/S3}
\input{build/tables/S4}

\clearpage
%=============================================================================
\section{Discriminative validity}
%=============================================================================

\input{build/tables/S5}

\clearpage
%=============================================================================
\section{Redundancy structure}
%=============================================================================

\input{build/tables/S6}
\input{build/tables/S7}
\input{build/tables/S8}
\input{build/tables/S9}

\clearpage
%=============================================================================
\section{Selection stability}
%=============================================================================

\input{build/tables/S10}

\begin{figure}[!htbp]
\centering
\renewcommand{\thefigure}{S2}
\includegraphics[width=\linewidth]{fig_S2_selection_frequency.pdf}
\caption{Selection frequency of each descriptor across the 50 subject-grouped
training partitions of the stability analysis, for the label-agnostic
reduced set and the three supervised selectors at matched cardinality. The
values are those of Supplementary Table~S10. On qEEG the label-agnostic
pipeline retains the same descriptors almost every time, which is why its
Nogueira index is the highest of the four (Table~14). On TCD it alternates
between members of the velocity family, whose correlations are nearly tied,
while ReliefF retains almost the same subset every time.}
\label{fig:S2}
\end{figure}

\clearpage
%=============================================================================
\section{Analysis of variance}
%=============================================================================

\input{build/tables/S11}

\FloatBarrier
\end{document}
